\documentclass{simple}

\title[How to read the room]{How to read the room}
\institute{Excursie, Bușteni}
\author[Razvan Deaconescu]{razvand@gmail.com}
\date{14 martie 2026}

\begin{document}

\frame{\titlepage}

\begin{frame}{Cine inițiază flirtul?}
  \centering
  \pause
  \Large
  el sau ea? \\
  \pause
  \vspace{1cm}
  De ce? \\
  \pause
  \vspace{1cm}
  \small
  \url{https://x.com/Devon_Eriksen_/status/1918506702267175234}
\end{frame}

\begin{frame}{Tipuri de comunicare}
  \centering
  \pause
  \Large
  verbal, para-verbal, non-verbal \\
  \pause
  \vspace{1cm}
  directă, indirectă (subtilă, cu subînțeles) \\
  \pause
  \vspace{1cm}
  text, online, în persoană
\end{frame}

\begin{frame}{Ce înseamnă?}
  \centering
  \pause
  \Large
  Da \\
  \pause
  \vspace{1cm}
  Ai dreptate \\
  \pause
  \vspace{1cm}
  Ce prost ești
\end{frame}

\begin{frame}{}
  \centering
  \pause
  \Large
  Comunicarea para-verbală / non-verbală e, în general, congruentă cu cea verbală. \\
  \pause
  \vspace{1cm}
  Dar se poate întâmplă să nu fie: ,,Her mouth said no, but her eyes said yes.''
\end{frame}

\begin{frame}{}
  \centering
  \pause
  \Large
  Comunicarea para-verbală / non-verbală este: \\
  \pause
  \vspace{1cm}
  mai onestă, mai greu de ,,mințit'' \\
  \pause
  \vspace{1cm}
  subtilă, cu subînțeles
\end{frame}

\begin{frame}{De ce subtilitate, de ce cu subînțeles?}
  \begin{itemize}
    \pause \item \textit{plausible deniability}
    \pause \item arătat cu degetul, jenat
    \pause \item comunicare selectivă (selectezi din audiență - cine ,,prinde'')
    \pause \item evaluarea cuiva, \textit{vetting} - se prinde?
    \pause \item nu vrei să fii nepoliticos, perceput ca bădăran, necioplit, agresiv, neceremonios
  \end{itemize}
\end{frame}

\begin{frame}{Important: Read the room}
  \pause Să ,,percepi atmosfera din încăpere''
  \begin{itemize}
    \pause \item să vezi starea de spirit / emoțională a oamenilor
    \pause \item să percepi nivelul de (dis)confort
    \pause \item să înțelegi cât de mult chef / dorință au de ceva cei din încăpere
    \pause \item să prinzi ,,semnele''
    \pause \item să (re)acționezi cât mai eficace ținând cont de atmosfera din cameră
  \end{itemize}
\end{frame}

\begin{frame}{Exemple}
  \begin{itemize}
    \item Amâni o discuție dacă vezi că persoana e ocupată sau nu are chef.
    \item Nu faci glume nepotrivite când e o stare de spirit foarte serioasă.
    \item Abordezi pe cineva când detectezi "semnele".
    \item Deschizi conversații când vezi că cineva are nevoie de sprijin.
    \item Oferi ajutor atunci când cineva are dificultăți.
  \end{itemize}
  \pause Toate au excepții care depind de context.
\end{frame}

\begin{frame}{De ce e dificil să percepi atmosfera?}
  \centering
  \pause
  \Large
  inner-world vs outer-world \\
  \pause
  \vspace{1cm}
  \textit{Your outer world is a reflection of your inner world.} \\
  \vspace{3mm}
  \hfill \textit{T. Harv Eker} \\
  \pause
  \vspace{1cm}
  experiență \\
  \pause
  \vspace{1cm}
  lipsă de timp / ,,lățime de bandă''
\end{frame}

\begin{frame}{Mindset}
  \begin{itemize}
    \pause \item \textit{get out of yourself}
    \pause \item plăcerea ascultatului, observatului
    \pause \item detașare / ,,dispassion'' (nu lua foarte personal, sau foarte în piept / frontal)
  \end{itemize}
\end{frame}

\begin{frame}{Bune practici}
  \pause Nu există rețete.
  \begin{itemize}
    \pause \item Redu consumul de energie.
      \begin{itemize}
        \pause \item Nu vorbi / exprima foarte mult.
        \pause \item Nu-ți încărca creierul cu foarte multe lucruri (în acele momente).
        \pause \item Mișcă-te mai puțin / mai eficient.
      \end{itemize}
    \pause \item Fă pauze, ia timp să asculți, să observi.
    \pause \item Redu stimulii: social media, muzică.
    \pause \item Fă o trecere prin experiențe, rememorează, analizează, judecă, creează scenarii alternative (constructiv, nu ca reproș).
    \pause \item Fii conștient de bias-urile, credințele și triggerele tale și încearcă să nu te afecteze în percepție sau în reacție.
  \end{itemize}
\end{frame}

\begin{frame}{How to make the room readable}
  \pause Comunicare asertivă:
  \begin{itemize}
    \pause \item francă, deschisă, completă, congruentă (verbal, para-verbal, non-verbal), constructivă, factuală
  \end{itemize}
  \pause Adică nu:
  \begin{itemize}
    \pause \item agresivă, pasiv-agresivă, defensivă, victmizantă, centrată pe sine, opinionated
  \end{itemize}
\end{frame}

\begin{frame}{De avut în vedere}
  \centering
  \pause
  \Large
  Emoțiile, conexiunile personale, intime merg mână în mână cu subtilitatea și subînțelesul. \\
  \pause
  \vspace{1cm}
  Nu faceți comunicare asertivă când e vorba de emoții, conexiuni personale.
\end{frame}

\begin{frame}{}
  \centering
  \pause
  \Large
  \textit{The quieter you become, the more you are able to hear.}
  \vspace{3mm}
  \hfill \textit{Jalāl al-Dīn Muḥammad Rūmī}
\end{frame}

\begin{frame}{Link la slide-uri}
  \url{https://www.slideshare.net/razvandeaconescu}
\end{frame}

\end{document}
